\subsection{Inverses in an Equation} 
Recall that we form the inverse of a relation by reversing the \makeblanksmall and \makeblanksmall coordinate of each pair. If the relation is given in set-builder notation using an equation, we can do this in one of two ways.
\begin{Example}
Let $R=\Set{(x,y)\mid x^2-y^2=9}$. Write its inverse in two ways.
\begin{enumerate}
\item We can reverse the $x$ and $y$ in the \makeblank.\lbr$R\inv=\Set{\makeblank}$
\item We can reverse the $x$ and $y$ in the \makeblank.\lbr$R\inv=\Set{\makeblank}$
\end{enumerate}
\end{Example}
\noindent
Since, when working with functions, we tend to think of $x$ as the \makeblank[1in] and $y$ as the \makeblank[1in], it makes more sense to keep the \makeblank[1in] and reverse $x$ and $y$ in the \makeblank.
\begin{Example}
Let $S=\Set{(x,y)\mid y=2x-6}$. Write $S\inv$, then solve the equation in $S\inv$ for $y$.\lbr
$S\inv=\Set{\makeblank}$
\begin{lpic}{}{5}\regtextleft{0}{-1}{Solving for $y$:};\end{lpic}
So $S\inv=\Set{\makeblank}$
\end{Example}